\documentclass[12pt]{article}

\usepackage[spanish]{babel}
\usepackage[utf8]{inputenc}
\usepackage[T1]{fontenc}
\usepackage[a4paper,margin=2.5cm]{geometry}
\usepackage{graphicx}
\usepackage{setspace}
\usepackage{titlesec}

\geometry{margin=1in}
\onehalfspacing

\title{Monitoreo y Diagnóstico de Redes}
\author{}
\date{}

\begin{document}

\begin{titlepage}
    \centering

    % Logo opcional
    %\includegraphics[width=0.3\textwidth]{logo.png}\par\vspace{1cm}

    {\scshape\Large Monitoreo y Diagnóstico de Redes \par}
    \vspace{1.5cm}

    {\huge\bfseries Evidencia del trabajo \par}
    \vspace{0.5cm}
    {\Large Tarea 2 \par}

    \vfill

    {\large
    Nombre del Alumno: Emilio Izquierdo Montero {\hspace{6cm}} \\
    \vspace{0.3cm}
    Profesor: Pedro Pablo Martinez Carmona {\hspace{6cm}} \\
    \vspace{0.3cm}
    Grupo / Semestre: 8 {\hspace{6cm}} \\
    \vspace{0.3cm}
    Fecha: 17 de abril del 2026 {\hspace{6cm}}
    \par}

    \vfill

\end{titlepage}
\section*{1. ¿Para qué pueden emplearse las herramientas de diagnóstico?}

Las herramientas de diagnóstico sirven para supervisar, analizar y detectar problemas en una red o sistema. Permiten medir el rendimiento, identificar fallos, monitorear procesos y verificar el estado de los dispositivos.

Algunas herramientas comunes son:

\begin{itemize}
    \item \textbf{htop}: monitoriza procesos y uso de CPU, memoria y carga del sistema.
    \item \textbf{top}: muestra procesos activos y consumo de recursos.
    \item \textbf{ping}: verifica conectividad y latencia.
    \item \textbf{traceroute}: muestra la ruta seguida por los paquetes.
    \item \textbf{nmap}: escaneo de puertos y análisis de hosts.
    \item \textbf{Wireshark}: captura y análisis de paquetes.
\end{itemize}

\section*{2. Ventaja de NetFlow}

\textbf{Respuesta correcta:}

\begin{itemize}
    \item b. ayuda a los análisis de seguridad
\end{itemize}

\section*{3. Verdadero o falso}

\subsection*{a)}
La monitorización activa implica el envío de traps por parte de agentes.

\textbf{Falso}

Los traps pertenecen a la monitorización pasiva mediante SNMP.

\subsection*{b)}
Nmap es una herramienta para monitorizar el rendimiento.

\textbf{Falso}

Nmap se utiliza principalmente para escaneo de puertos y descubrimiento de hosts.

\section*{4. ¿Cuáles son las tres fases de trabajo de Cacti?}

\begin{enumerate}
    \item Recolección de datos.
    \item Almacenamiento de información.
    \item Generación de gráficas.
\end{enumerate}

\section*{5. ¿Para que NetFlow funcione es necesario tenerlo habilitado en todos los routers de la red?}

\textbf{Falso}

No es obligatorio habilitar NetFlow en todos los routers. Puede activarse únicamente en los dispositivos clave donde se desea analizar el tráfico de red.

\section*{6. ¿Qué hace la herramienta traceroute?}

Traceroute permite identificar la ruta que siguen los paquetes desde un origen hasta un destino, mostrando los saltos intermedios y el tiempo de respuesta de cada uno.

\section*{7. Con MRTG se puede hacer...}

\textbf{Respuesta correcta:}

\begin{itemize}
    \item c. un gráfico que muestre el tráfico de los dispositivos administrados
\end{itemize}

\section*{8. Explique qué tipo de pruebas se hacen para controlar la monitorización}

Se realizan pruebas de conectividad, latencia, pérdida de paquetes, consumo de ancho de banda, uso de CPU, memoria y disponibilidad de servicios.

\section*{9. ¿Para usar Nagios es necesario un servidor web?}

\textbf{Sí}

Nagios utiliza una interfaz web para visualizar alertas, estados y reportes, por lo que normalmente requiere un servidor web como Apache.

\section*{10. Indique dos elementos que conforman Cricket}

\begin{itemize}
    \item RRDTool
    \item SNMP
\end{itemize}

\section*{11. Indique un cliente de Nagios}

\textbf{Respuesta correcta:}

\begin{itemize}
    \item d. Nsclient++
\end{itemize}

\section*{12. Para elaborar un gráfico con Cacti...}

\textbf{Respuesta correcta:}

\begin{itemize}
    \item c. se puede escoger el OID de la MIB
\end{itemize}

\section*{13. Relacione cada elemento}

\begin{center}
\begin{tabular}{|l|l|}
\hline
Zenmap & Testeo de puertos \\
\hline
Nagios & Nrpe \\
\hline
Netflow & Cisco \\
\hline
J-Flow & Juniper \\
\hline
\end{tabular}
\end{center}

\section*{14. Protocolo de administración de red}

\textbf{Respuesta correcta:}

\begin{itemize}
    \item d. RMON
\end{itemize}

\section*{15. ¿Qué conceptos deben considerarse en la planificación del análisis de rendimiento?}

\begin{itemize}
    \item Latencia
    \item Jitter
    \item Pérdida de paquetes
    \item Uso de ancho de banda
    \item Disponibilidad
\end{itemize}

\section*{16. Para valorar la utilización del canal se puede emplear...}

\textbf{Respuesta correcta:}

\begin{itemize}
    \item c. el 95-percentil
\end{itemize}

\section*{17. Verdadero o falso}

\subsection*{a)}
Intentar mitigar el Jitter puede aumentar la latencia.

\textbf{Verdadero}

\subsection*{b)}
Entre los destinatarios de la información de los análisis pueden estar los clientes de la empresa.

\textbf{Verdadero}

\section*{18. ¿El clustering mejora el rendimiento de los servidores?}

\textbf{Verdadero}

El clustering distribuye la carga entre varios servidores y mejora disponibilidad y rendimiento.

\section*{19. Diferencia entre indicador y métrica}

\begin{itemize}
    \item \textbf{Indicador}: representa el estado general de un sistema.
    \item \textbf{Métrica}: valor numérico específico medido.
\end{itemize}

\section*{20. ¿Qué soluciones existen para garantizar disponibilidad?}

\begin{itemize}
    \item Redundancia
    \item Clustering
    \item Balanceo de carga
    \item Sistemas UPS
    \item Respaldos
\end{itemize}

\section*{21. Elemento usado para calcular el retardo de transmisión}

\textbf{Respuesta correcta:}

\begin{itemize}
    \item c. el tamaño del paquete
\end{itemize}

\section*{22. ¿Qué se puede utilizar para medir el retardo?}

\textbf{Respuesta correcta:}

\begin{itemize}
    \item d. Todas las respuestas anteriores son correctas
\end{itemize}

\section*{23. ¿Qué se hace para evitar la variación del retardo?}

Se implementan mecanismos de Calidad de Servicio (QoS), priorización de tráfico y control de congestión.

\section*{24. ¿Los discos duros son los dispositivos más importantes de monitorizar?}

\textbf{Sí}

Porque almacenan información crítica y su fallo puede afectar el funcionamiento completo del sistema.

\section*{25. Indique tres elementos que limitan la capacidad efectiva de un canal}

\begin{itemize}
    \item Congestión
    \item Latencia
    \item Pérdida de paquetes
\end{itemize}

\section*{26. Principal causa de pérdida de paquetes}

\textbf{Respuesta correcta:}

\begin{itemize}
    \item c. la congestión de la red
\end{itemize}

\section*{27. Elemento usado para calcular el retardo extremo a extremo}

\textbf{Respuesta correcta:}

\begin{itemize}
    \item a. el retardo de propagación
\end{itemize}

\section*{28. Relacione cada elemento}

\begin{center}
\begin{tabular}{|l|l|}
\hline
Retardo procesamiento & Capacidad de CPU \\
\hline
Jitter & Variación del retardo \\
\hline
Retardo propagación & Velocidad de la luz \\
\hline
Retardo transmisión & Tasa de envío \\
\hline
\end{tabular}
\end{center}

\section*{29. Entre las medidas correctivas se encuentra...}

\textbf{Respuesta correcta:}

\begin{itemize}
    \item d. modificar el uso de la red
\end{itemize}

\end{document}
